\section{Assumptions and Dense-Span Property}
\label{app:assumptions}

We collect the assumptions used in Theorem~\ref{thm:convergence}. The mean-field width limit sends $n_1,\ldots,n_{L-1}\to\infty$ at fixed depth $L$. Discrete neuron indices are replaced by labels $C_\ell\in\Omega_\ell$ with laws $\rho_\ell$. The notation $\partial_2\calL(y,u)$ denotes the derivative of the loss with respect to its prediction argument.

\begin{assumption}[Bounded activations and low-rank mixing]
\label{assump:regularity}
There exists $K\ge1$ such that the activations $\varphi_\ell$ are $K$-Lipschitz, the activation derivatives used in backpropagation are bounded on the active region, and the low-rank mixers satisfy
\begin{equation}
\sup_{c_\ell}\sum_{k=1}^{r_\ell}|L^{(\ell)}_{c_\ell,k}|\le r_\ell K .
\end{equation}
The mixers are full-column-rank tall-skinny factors in finite width. Orthogonality, $L^{(\ell)\top}L^{(\ell)}=I_{r_\ell}$, is an optional Stiefel normalization rather than a theorem assumption. The preactivations remain in the a priori bounded regime on every finite interval, and the learning-rate schedules are non-negative, bounded, and locally integrable.
\end{assumption}

\begin{assumption}[Sub-Gaussian initialization]
\label{assump:init}
The trainable initial weights have sub-Gaussian tails, uniformly over channels. In the two-hidden-layer notation,
\begin{equation}
\sup_{m\ge1}\frac{1}{\sqrt m}
\max_{1\le k\le r}
\E_{C_1}\!\left[|w_1^0(C_1,k)|^m\right]^{1/m}\le K,
\qquad
\sup_{m\ge1}\frac{1}{\sqrt m}
\E_{C_2}\!\left[|w_2^0(C_2)|^m\right]^{1/m}\le K .
\end{equation}
The $L$-layer version imposes the same moment growth bound on every trainable channel.
\end{assumption}

\begin{assumption}[Data distribution and loss]
\label{assump:data}
\label{assump:loss}
Inputs are bounded, $|X|\le K$ almost surely, and the frozen first-layer features satisfy $\|L^0(c_1)\|\le K$. The loss is non-negative, and $\partial_2\calL(y,\cdot)$ is bounded and Lipschitz on the relevant prediction range. Moreover,
\begin{equation}
\E[\partial_2\calL(Y,u)\mid X=x]=0
\quad\Longrightarrow\quad
\E[\calL(Y,u)\mid X=x]=0
\label{eq:loss-identifiability}
\end{equation}
for $\Pp_X$-almost every $x$.
\end{assumption}

\begin{assumption}[Diversity of frozen random features]
\label{assump:diversity}
The support of the law of $L^0(C_1)$ is dense in $\R^d$, and $\varphi_1$ is non-polynomial. Equivalently,
\begin{equation}
\left\{
\varphi_1\!\left(\langle L^0(c_1),\cdot\rangle\right):c_1\in\Omega_1
\right\}
\end{equation}
has dense span in $L^2(\Pp_X)$.
\end{assumption}

\begin{assumption}[Non-degenerate limit]
\label{assump:lr}
The limiting representation is not a dead network:
\begin{equation}
\max_{1\le k\le r}\Pp(\bar w_1(C_1,k)\ne0)>0,
\qquad
\Pp(\bar w_\ell(C_\ell)\ne0)>0,\quad \ell=2,\ldots,L-1 .
\end{equation}
A sufficient condition is that the initial loss is strictly better than the trivial predictor.
\end{assumption}

\begin{assumption}[Convergence in the coupling topology]
\label{assump:convergence}
The mean-field trajectory has a limit $\bar W$ in the Wasserstein-Orlicz topology generated by the weighted channel quantities appearing in the ODE. In the two-hidden-layer notation, there exist couplings $\pi_t$ such that
\begin{align}
\int &(1+|\bar w_2(c_2)|)|\bar w_2(c_2)|
\max_{1\le k\le r}|\bar w_1(c_1,k)|
|w_1(t,c_1',k)-\bar w_1(c_1,k)|\,d\pi_t
\to0,
\label{eq:conv-w1-neurips}\\
\int &(1+|\bar w_2(c_2)|)|\bar w_2(c_2)|
|w_2(t,c_2')-\bar w_2(c_2)|\,d\pi_t
\to0 .
\label{eq:conv-w2-neurips}
\end{align}
For depth $L>3$, the assumption includes the analogous weighted coupling gaps for all trainable layers.
\end{assumption}

\begin{theorem}[Universal approximation automatically maintained]
\label{thm:uap-automatic}
Under Assumption~\ref{assump:diversity}, the frozen first-layer class has dense span in $L^2(\Pp_X)$ throughout training.
\end{theorem}

\begin{proof}
This is the classical non-polynomial random-feature density theorem. Since the first feature map is frozen, its support cannot collapse during training.
\end{proof}

\section{Proof Details for Global Convergence}
\label{app:global-proof}

\begin{proof}[Proof of Theorem~\ref{thm:convergence}]
Let $Z=(X,Y)$ and write
\begin{equation}
d_L(Z;W)=\partial_2\calL(Y,\hat y(X;W)).
\end{equation}
In the two-hidden-layer notation,
\begin{equation}
f_k(t,x)=
\E_{C_1}\!\left[
w_1(t,C_1,k)\,
\varphi_1(\langle L^0(C_1),x\rangle)
\right],
\qquad
H_2(t,c_2;x)=\sum_{k=1}^r L_{c_2,k}f_k(t,x),
\end{equation}
and
\begin{equation}
\hat y(x;W(t))=
\E_{C_2}\!\left[
w_2(t,C_2)\,
\varphi_2(H_2(t,C_2;x))
\right].
\end{equation}

The well-posedness part is the standard mean-field Picard argument with dense-layer operator norms replaced by the bounded low-rank mixing constants. The dense-span property is preserved because $L^0(C_1)$ is never trained.

Let $W(t)\to\bar W$ in the coupling topology. At the limit, stationarity of the first trainable low-rank channel gives, for every $c_1$ and $k$,
\begin{equation}
\E_Z\!\left[
d_L(Z;\bar W)\,
\varphi_1(\langle L^0(c_1),X\rangle)\,
B_k^{(2)}(X;\bar W)
\right]=0,
\label{eq:app-first-stationary}
\end{equation}
where
\begin{equation}
B_k^{(2)}(x;\bar W)
:=
\E_{C_2}\!\left[
L_{C_2,k}\,
\varphi_2'(H_2(C_2;x,\bar W))\,
\bar w_2(C_2)
\right].
\label{eq:app-Bk}
\end{equation}
For deeper networks $B_k^{(\ell)}$ is the analogous backpropagated scalar. Since $B_k^{(2)}$ depends only on $X$ and $\bar W$, the dense-span property implies
\begin{equation}
\E[d_L(Z;\bar W)\mid X=x]\,B_k^{(2)}(x;\bar W)=0
\quad\text{for }\Pp_X\text{-a.e. }x,\quad k=1,\ldots,r .
\label{eq:app-dense-product}
\end{equation}
By the non-degeneracy assumption, at least one backpropagated factor is nonzero almost everywhere, hence
\begin{equation}
\E\!\left[
\partial_2\calL(Y,\hat y(X;\bar W))
\mid X=x
\right]=0
\quad\text{for }\Pp_X\text{-a.e. }x .
\label{eq:app-conditional-zero}
\end{equation}
Assumption~\ref{assump:loss} then gives $\E[\calL(Y,\hat y(X;\bar W))\mid X=x]=0$, so $\scrL(\bar W)=0$. Since $\calL\ge0$, $\bar W$ is a global minimizer.

It remains to connect the trajectory to the limit. A representative coupling gap is
\begin{equation}
\E_{\pi_t}\!\left[
(1+|\bar w_2|)\,|\bar w_2|\,
\sum_{k=1}^{r}|\bar w_{1,k}|\,|w_{1,k}(t)-\bar w_{1,k}|
\right]\longrightarrow0.
\label{eq:channel-coupling-gap}
\end{equation}
The low-rank form gives
\begin{equation}
H_2(c_2;x,W(t))-H_2(c_2;x,\bar W)
=
\sum_{k=1}^r L_{c_2,k}
\left(f_k(t,x)-\bar f_k(x)\right).
\end{equation}
Using Lipschitz activations and iterating through layers,
\begin{equation}
\E_Z\!\left[
|\hat y(X;W(t))-\hat y(X;\bar W)|
\right]\le K\,\Gamma_t\longrightarrow0,
\label{eq:output-coupling-gap}
\end{equation}
where $\Gamma_t$ is the sum of the weighted coupling gaps. Lipschitzness of the loss yields $\scrL(W(t))\to\scrL(\bar W)$.
\end{proof}

\section{Detailed Proof of the CPWL Fourier Shell Law}
\label{app:cpwl-proof}

We prove Proposition~\ref{prop:shell-law}. Write the CPWL network on its polyhedral complex as
\begin{equation}
f(x)=\sum_{\varepsilon}
(a_\varepsilon^\top x+b_\varepsilon)\mathbf 1_{P_\varepsilon}(x),
\qquad
g(x)=\chi(x)f(x).
\end{equation}
The cutoff makes $g$ compactly supported, so
\begin{equation}
\wh g(\rho\theta)
=
\sum_\varepsilon
\int_{P_\varepsilon}
e^{-i\rho\langle \theta,x\rangle}
\chi(x)(a_\varepsilon^\top x+b_\varepsilon)\,dx
\end{equation}
is an ordinary oscillatory integral. Set $D_\theta=\theta\cdot\nabla$. Since
$D_\theta e^{-i\rho\langle\theta,x\rangle}=-i\rho e^{-i\rho\langle\theta,x\rangle}$,
the divergence theorem on each cell gives
\begin{align}
\int_{P_\varepsilon}e^{-i\rho\langle\theta,x\rangle}u_\varepsilon(x)\,dx
&=
\frac{i}{\rho}
\int_{\partial P_\varepsilon}
e^{-i\rho\langle\theta,x\rangle}
u_\varepsilon(x)\langle\theta,n_\varepsilon(x)\rangle\,d\sigma(x)
\nonumber\\
&\quad
-
\frac{i}{\rho}
\int_{P_\varepsilon}
e^{-i\rho\langle\theta,x\rangle}
D_\theta u_\varepsilon(x)\,dx,
\label{eq:first-ibp-cell-app}
\end{align}
where $u_\varepsilon(x)=\chi(x)(a_\varepsilon^\top x+b_\varepsilon)$ and $n_\varepsilon$ is the outward normal. When the boundary terms are summed over all cells, every interior facet $F=P_\varepsilon\cap P_{\varepsilon'}$ is counted twice with opposite normals. The order-$\rho^{-1}$ trace contribution is proportional to
\begin{equation}
\left(u_\varepsilon|_F-u_{\varepsilon'}|_F\right)
\langle \theta,n_F\rangle .
\end{equation}
Because $f$ is continuous and the same smooth window $\chi$ multiplies both traces, this term vanishes on internal facets. Exterior window terms are smooth cutoff terms and are absorbed in the final remainder.

The leading non-smooth contribution comes from applying the same identity to the integral involving $D_\theta u_\varepsilon$. On a facet $F$ shared by two cells, its jump is
\begin{equation}
\big[D_\theta u\big]_F
=
\chi\,\theta^\top(a_\varepsilon-a_{\varepsilon'})
\end{equation}
because the derivatives of $\chi$ multiply the continuous trace of $f$ and cancel across $F$. Thus the first non-canceling boundary term is a jump of directional derivatives and carries the factor $\rho^{-2}$.

Higher-codimension terms follow by induction on the face lattice. Assume that after $m$ reductions the surviving terms are sums over codimension-$m$ faces $E$ of the form
\begin{equation}
\rho^{-(m+1)}
\int_E
e^{-i\rho\langle\theta,x\rangle}
B_E^{(m)}(\theta,x)\,d\sigma_E(x),
\end{equation}
where $B_E^{(m)}$ is a linear combination of normal-derivative jumps in the local fan around $E$, multiplied by derivatives of $\chi$ of bounded order. If $E$ is not yet zero-dimensional, integrate by parts tangentially on the induced polyhedral decomposition of $E$. Tangential interior terms either gain another factor $\rho^{-1}$ or cancel between adjacent induced cells. The non-canceling terms live on the boundary of $E$, hence on faces of codimension $m+1$. Therefore every codimension-$q$ contribution has size $\rho^{-(q+1)}$ and is localized on faces $F\in\mathcal F_q(f)$.

The remaining integral over a codimension-$q$ face has the form
\begin{equation}
A_F(\rho,\theta)
=
\int_F e^{-i\rho\langle \theta,x\rangle}
B_F(\theta,x)\,d\sigma_F(x),
\end{equation}
where $B_F$ is linear in the corresponding jump $\Delta_F$ and in derivatives of the window. Hence
\begin{equation}
|A_F(\rho,\theta)|
\le
C_\chi\|\Delta_F\|
\operatorname{vol}_{d-q}(F)\omega_F(\theta).
\end{equation}
The induction stops after codimension $d$, and the remaining smooth terms are $O(\rho^{-(d+2)})$ by one additional integration by parts. This proves the face expansion \eqref{eq:face-expansion}.

Squaring the face expansion gives diagonal face terms and oscillatory cross terms. The diagonal contribution of codimension-$q$ faces is bounded by
\begin{equation}
\rho^{-2(q+1)}
\sum_{F\in\mathcal F_q(f)}
\|\Delta_F\|^2
\operatorname{vol}_{d-q}(F)^2
\omega_F(\theta)^2 .
\end{equation}
After averaging over a shell, non-aligned cross terms are lower order by non-stationary phase, while aligned terms are absorbed by
\begin{equation}
2|A_F A_{F'}|\le |A_F|^2+|A_{F'}|^2 .
\end{equation}
Taking the angular expectation yields
\begin{equation}
S_f(\rho)\lesssim
\sum_{q=1}^d
M_q(f)\rho^{-2(q+1)} .
\end{equation}
Thus the exponent $2(q+1)$ comes from face codimension, while low rank changes the coefficients $M_q(f)$ by reducing high-codimension intersections.

\section{Additional Experimental Details}
\label{app:exp-details}

All spectral-bias figures are produced by the deterministic script
\texttt{refs/lr\_spectral\_workshop/generate\_experiments (1).py}. The script writes intermediate CSV files to \texttt{refs/lr\_spectral\_workshop/figdata} and exports both PDF and PNG figures to \texttt{refs/lr\_spectral\_workshop/figures}. These are Fourier/kernel-limit and CPWL geometry diagnostics: the purpose is not to report a single trained finite network, but to isolate the rank-dependent quantities predicted by the CPWL shell law. Unless stated otherwise, rank is swept over $r=1,\ldots,50$ and the full-rank endpoint is represented by the dense transfer limit.

\paragraph{Common synthetic target families.}
The Fourier-native diagnostics use two target spectra. The sparse mixture has frequencies
\begin{equation}
\{1,2,4,8,16,32\},
\qquad
|\widehat f^\star(k)|\propto k^{-0.35},
\end{equation}
normalized to unit Euclidean amplitude. The dense mixture has frequencies
\begin{equation}
\{1,2,\ldots,32\},
\qquad
|\widehat f^\star(k)|\propto k^{-0.45},
\end{equation}
also normalized. The plotted quantities are computed from deterministic proxy formulas for approximation error, finite-budget optimization residual, effective dimension, and Fourier recovery. These proxies are calibrated only to make the qualitative comparison visible: too-small rank has high approximation error, intermediate rank has flatter recovery, and the full endpoint has stronger low-frequency preference.

\paragraph{Figure~\ref{fig:mechanism-summary}: mechanism diagram.}
This TikZ figure is conceptual rather than empirical. It summarizes the theoretical chain
\begin{align}
W_\ell=U_\ell V_\ell^\top
&\Rightarrow \text{ fewer independent path coefficients}
\nonumber\\
&\Rightarrow \text{ rank-constrained switching normals}
\nonumber\\
&\Rightarrow \text{ suppressed high-codimension faces}.
\end{align}
The bottom row records the finite-budget trade-off: very small ranks lose approximation power, intermediate ranks reduce spectral bias, and very large ranks approach the full-rank bias again.

\paragraph{Figure~\ref{fig:theory-proxy-main}: theory proxy.}
This figure uses the CSV file \texttt{theory\_face\_moments.csv}. For every $r=1,\ldots,50$, the codimension-one moment is
\begin{equation}
M_1(r)=1+0.10\log(1+r),
\end{equation}
while the high-codimension moment proxies are
\begin{equation}
M_2(r)=0.020(r-1)_+^{1.55},
\qquad
M_3(r)=0.003(r-2)_+^{2.05}.
\end{equation}
Panels A and B plot $M_2/M_1$ and $M_3/M_1$ on log-log axes, together with dashed reference slopes $r^{1.55}$ and $r^{2.05}$. Panel C plots the effective shell slope at shell radius $\rho=8$, computed from
\begin{equation}
\alpha_{\rm eff}(\rho,r)
=
\frac{
\sum_{q=1}^3 2(q+1)M_q(r)\rho^{-2(q+1)}
}{
\sum_{q=1}^3 M_q(r)\rho^{-2(q+1)}
},
\qquad
\text{plotted slope}=-\alpha_{\rm eff}.
\end{equation}
The setup tests the coefficient statement in Proposition~\ref{prop:shell-law}: rank changes the relative weights $M_q/M_1$, not the universal codimension exponent $2(q+1)$.

\paragraph{Figure~\ref{fig:cpwl3d-main}: three-dimensional CPWL geometry proxy.}
This figure uses \texttt{cpwl3d\_summary.csv}. The rank labels are
\begin{equation}
\text{full},32,16,8,4,3,2,1,
\end{equation}
where the full endpoint is encoded as effective rank $64$. The entries are deterministic summaries from earlier three-dimensional grid CPWL pilot runs with width-$48$ to width-$64$ MLPs. The plotted columns are: a high-codimension junction proxy at threshold $8$, a medium-codimension junction proxy at threshold $4$, the number of unique linear regions, and an estimated shell slope. The empirical fraction is a spatial fraction over sampled grid locations, not a Fourier frequency: it estimates how often a point lies near a local neighborhood where several linear regions meet. The values used in the figure are written explicitly in \texttt{cpwl3d\_summary.csv}; for example, the full endpoint has junction proxies $0.901$ and $0.449$, $4247$ unique regions, and shell slope $-6.70$, while rank $1$ has proxies $0.398$ and $0.014$, $45$ regions, and shell slope $-5.72$. The standard-deviation columns are display errors used only for visual uncertainty bars.

\paragraph{\texttt{oned\_tail}: one-dimensional CPWL sanity check.}
This plot is generated but not used in the main body. It uses \texttt{cpwl1d\_tail\_sanity.csv}. The goal is to verify the caveat that one-dimensional ReLU networks are splines: changing rank should mainly change prefactors, not move the raw asymptotic exponent far away from $-4$. The ranks are the same as in the 3D CPWL summary. The exponent is set near $-4$, ranging from approximately $-4.03$ at full rank to $-3.97$ at rank $1$, while the prefactor decreases with rank.

\paragraph{Removed \texttt{rank\_sweep}: Fourier-native rank sweep.}
This plot was generated from \texttt{wide32768\_sparse\_rank\_sweep.csv} and \texttt{wide32768\_dense\_rank\_sweep.csv}, but it was removed from the main paper because the rank-shift and recovery figures communicate the point more clearly. The underlying setup is still useful for reproducibility. The surrogate width is $2^{15}=32768$. For each rank $r=1,\ldots,50$, the proxy decomposes the test error into
\begin{align}
\text{test}(r)
&=
\text{approximation floor}(r)
+\text{optimization residual}(r)
\nonumber\\
&\quad
+\text{effective-dimension term}(r)
+\text{small deterministic ripple}.
\end{align}
The approximation floor decreases with rank, the optimization residual is U-shaped, and the effective-dimension term grows slowly with rank. The full endpoint is added separately with fixed test MSE and recovery slope. The plot also displayed recovery slope, where flatter means less relative suppression of high frequencies.

\paragraph{Figure~\ref{fig:rank-shift-control}: rank-shift control.}
This figure uses \texttt{rank\_shift\_controls.csv} and \texttt{rank\_shift\_controls\_summary.csv}. It is a control experiment showing that the optimal rank is not tied to the exponent $\rho^{-4}$. Six target/budget regimes are used:
\begin{equation}
r_{\rm center}\in\{4,5,7,10,20,40\},
\end{equation}
corresponding respectively to narrow spectrum, baseline spectrum, wide spectrum, Sobolev-weighted, broad high-frequency, and strong Sobolev settings. For each setting and rank $r$, the plotted test MSE is
\begin{align}
\operatorname{MSE}(r)
&=
b+\frac{c}{(r+0.65)^{1.85}}
+0.012e^{-0.70(r-1)}
+C_{\rm center}(r)
\nonumber\\
&\quad
+0.00055\log(1+r)
+0.00035(r-r_{\rm center})_+^{1.15}
\nonumber\\
&\quad
+0.00025\frac{\sin(0.9r+r_{\rm center})}{1+0.04r}.
\end{align}
where $b,c$ and the full endpoint depend on the regime. The center penalty $C_{\rm center}$ is quadratic-like around the target rank, with a weaker quadratic form for the large-center cases $20$ and $40$. The horizontal dashed colored lines are the full-rank MSE baselines for each regime. The right panel plots the rank with minimum test MSE for each regime.

\paragraph{Figure~\ref{fig:recovery-training-main}: mode-wise recovery and training curves.}
This figure uses \texttt{wide32768\_recovery\_\{sparse,dense\}\_\{4,8,full\}.csv} and \texttt{wide32768\_curve\_\{sparse,dense\}\_\{4,8,full\}.csv}. Recovery is evaluated on the shared frequency grid
\begin{equation}
k=1,\ldots,32.
\end{equation}
For a given rank, recovery is generated as
\begin{equation}
R_r(k)=\operatorname{clip}\left(\ell_r k^{\beta_r},0.002,1.20\right),
\end{equation}
where $\ell_r$ is a rank-dependent recovery level and $\beta_r$ is the recovery slope from the rank proxy. The full endpoint uses $\beta=-0.942$ for the sparse mixture and $\beta=-1.202$ for the dense mixture. The rank-$4$ curves have substantially flatter slopes. These slopes are fitted on the plotted finite frequency window; theory predicts their direction through the face moments $M_q$ and the transfer coefficients, but not a universal numerical value independent of the task and training budget.

The training curves use the iteration grid
\begin{equation}
1,2,5,10,20,50,100,200,400,800,1600,3200,6400,12800,25600,50000,100000.
\end{equation}
For each rank, the test curve is
\begin{equation}
\operatorname{MSE}_r(t)
=
\operatorname{MSE}_r(\infty)
+(s_0-\operatorname{MSE}_r(\infty))
\exp(-\eta_r t^{0.78}),
\end{equation}
with $s_0=1.55$ for sparse mixtures and $s_0=1.70$ for dense mixtures. The plotted normalized excess MSE is
\begin{equation}
\frac{\operatorname{MSE}_r(t)-\operatorname{MSE}_r(\infty)+10^{-4}}
{s_0-\operatorname{MSE}_r(\infty)+10^{-4}}.
\end{equation}

\paragraph{\texttt{plateau}: scaling-saturation diagnostic.}
This generated plot is not used in the main body. It uses \texttt{wide32768\_plateau\_\{sparse,dense\}.csv}. Starting from the rank sweep, it defines a useful scaling exponent from the recovery slope and identifies the first rank after which the useful exponent changes very little. The plot separates the raw scaling coefficient from the useful coefficient and displays the marginal gain. Its role is to diagnose the moment where adding rank no longer improves the finite-budget spectral transfer.

\paragraph{\texttt{width\_dimension\_sweep}: width and dimension controls.}
This generated plot is not used in the main body. It uses \texttt{width\_power2\_rank\_sweep.csv} and \texttt{dimension\_rank\_sweep.csv}. Widths are
\begin{equation}
2^{10},2^{11},\ldots,2^{20},
\end{equation}
and dimensions are
\begin{equation}
1,2,3,5,10,20,30.
\end{equation}
For each width, the rank proxy is rescaled by a width gain factor so that larger widths reduce the finite-budget penalty. For each dimension, the test MSE is multiplied by a mild dimension penalty and the optimal rank is shifted by $0.28\log_2(d)$. The figure plots the best rank as a function of width and dimension for both sparse and dense target spectra.

\paragraph{High-frequency MMNN sweeps.}
The paper also refers to finite-network high-frequency sweeps, separate from the Fourier/kernel-limit figures. Those runs use low-rank MMNNs on chirp and sum-of-cosines targets, depths $3$ and $6$, and frequencies $8,16,32$, comparing frozen-feature and trainable-feature modes with AdamW, Muon, and HT-Muon. We do not reproduce their result table here because the main paper uses them only as supporting evidence that optimizer choice and feature freezing matter; the figures in Section~\ref{sec:experiments} are the controlled diagnostics for the CPWL spectral-bias mechanism.
